% ═══════════════════════════════════════════════════════════════
%  chapters/ch16-debug.tex
%  Отладка: инструменты, типичные ошибки, приёмы
% ═══════════════════════════════════════════════════════════════

\chapter{Отладка}
\label{ch:debug}

\epigraph{<<Отладка~--- это в два раза сложнее, чем написание кода.
Если вы пишете код настолько умно, насколько можете, вы по
определению недостаточно умны, чтобы его отлаживать>>.}{---
Брайан Керниган}

\section{Введение}
\label{sec:debug-intro}

Отладка~--- это навык. Не <<магия>>, не <<талант>>~---
навык, который тренируется. Чем больше вы
отлаживаете, тем быстрее находите проблемы.

Эта глава~--- про инструменты отладки и типичные
ошибки в нашем приложении.

\begin{ideabox}
Отладка~--- это \emph{систематический} процесс, а не
\emph{случайный} поиск. Вы \emph{формулируете} гипотезу:
<<программа падает, потому что $X$>>. Потом \emph{проверяете}.
Если гипотеза \emph{неверна}~--- \emph{другая} гипотеза.
И так, пока не найдёте.
\end{ideabox}

\subsection{Что вы узнаете}

\begin{itemize}
    \item Как работает \texttt{debug\_log}.
    \item Как использовать \texttt{gdb}.
    \item Как найти утечки через \texttt{valgrind}.
    \item Что такое \texttt{ASan} и \texttt{UBSan}.
    \item Как читать предупреждения \texttt{GLib}.
    \item Типичные ошибки нашего приложения.
\end{itemize}

\section{Инструменты отладки}
\label{sec:debug-tools}

В Linux есть набор инструментов для отладки.
Каждый решает свою задачу. Выбор зависит
от типа проблемы.

\begin{figure}[htbp]
\centering
\adjustbox{max width=\textwidth}{%
\begin{tikzpicture}[
    x=1cm, y=1cm,
    problem/.style={
        draw, rounded corners=2mm, thick,
        minimum width=3.5cm, minimum height=1cm,
        align=center, font=\small,
        fill=clrWarn!20
    },
    tool/.style={
        draw, rounded corners=2mm, thick,
        minimum width=3cm, minimum height=0.8cm,
        align=center, font=\small\ttfamily,
        fill=clrLogicData!20
    },
    arr/.style={-{Latex[length=3mm]}, thick}
]

% ─── Координатная сетка (для отладки) ────────────
\draw[step=1cm, gray!20, very thin] (-6, -6) grid (6, 1);

\node[problem] (p1) at (-4, 0)   {Программа падает\\(segfault)};
\node[problem] (p2) at (-4, -1.5) {Утечка памяти};
\node[problem] (p3) at (-4, -3)   {Странное поведение};
\node[problem] (p4) at (-4, -4.5) {<<Плавающий>> баг};

\node[tool] (t1) at (3, 0)    {gdb};
\node[tool] (t2) at (3, -1.5) {valgrind};
\node[tool] (t3) at (3, -3)   {debug\_log};
\node[tool] (t4) at (3, -4.5) {ASan / UBSan};

\draw[arr] (p1) -- (t1);
\draw[arr] (p2) -- (t2);
\draw[arr] (p3) -- (t3);
\draw[arr] (p4) -- (t4);

% Метки осей (для отладки)
\foreach \x in {-6,...,6} {
    \node[font=\tiny, gray] at (\x, -6.3) {\x};
}
\foreach \y in {-6,...,1} {
    \node[font=\tiny, gray] at (-6.3, \y) {\y};
}

\end{tikzpicture}
}% end adjustbox
\caption{Инструменты отладки и их задачи.}
\label{fig:debug-tools}
\end{figure}

\section{Отладочный лог: \texttt{debug\_log.c}}
\label{sec:debug-log}

Первый инструмент — самый простой.
Это лог-файл, куда приложение пишет
все важные события. Работает всегда~---
даже когда нет gdb под рукой.

\begin{listing}[htbp]
\begin{minted}[fontsize=\footnotesize, frame=none, bgcolor=gray!5]{c}
#define DEBUG_LOG_PATH "/tmp/arduino_debug.log"

void debug_log_init(void)
{
    FILE *f = fopen(DEBUG_LOG_PATH, "w");
    if (f) {
        fclose(f);
    }
}

void debug_log(const char *format, ...)
{
    FILE *log_file = fopen(DEBUG_LOG_PATH, "a");
    if (!log_file) {
        return;
    }

    time_t rawtime = time(NULL);
    struct tm timeinfo;
    localtime_r(&rawtime, &timeinfo);

    fprintf(log_file, "[%02d:%02d:%02d] ",
            timeinfo.tm_hour, timeinfo.tm_min, timeinfo.tm_sec);

    va_list args;
    va_start(args, format);
    vfprintf(log_file, format, args);
    va_end(args);

    fputc('\n', log_file);
    fclose(log_file);
}
\end{minted}
\caption{\texttt{debug\_log.c}~--- отладочный лог}
\label{lst:debug-log}
\end{listing}

\subsection{Как работает}

\textbf{Три ключевых решения:}

\begin{enumerate}
    \item \textbf{Путь в \texttt{/tmp}}~--- не \texttt{\$HOME}.
          Лог~--- временный, его можно удалить.
    \item \textbf{Открытие на каждое сообщение}~---
          просто и надёжно. Не нужен mutex,
          не нужна обработка ротации.
    \item \textbf{Временные метки}~--- \texttt{HH:MM:SS}.
          Позволяют понять, когда произошло событие.
\end{enumerate}

\begin{warnbox}
\texttt{debug\_log} не потокобезопасен. Он открывает
файл на каждое сообщение. Если два потока
пишут одновременно~--- строки перемешаются.
В нашем приложении это не проблема~---
всё в главном потоке GTK.
\end{warnbox}

\subsection{Как использовать}

В коде~--- через макрос \texttt{LOG\_DEBUG}:

\begin{listing}[htbp]
\begin{minted}[frame=none, bgcolor=gray!5]{c}
LOG_DEBUG("Порт: получено %zu байт", bytes_read);
LOG_DEBUG("Открытие порта %s", port);
LOG_DEBUG("Автореконнект запущен (интервал %d мс)",
          RECONNECT_INTERVAL_MS);
\end{minted}
\caption{Использование \texttt{LOG\_DEBUG}}
\label{lst:debug-log-usage}
\end{listing}

В терминале~--- наблюдение в реальном
времени:

\begin{minted}[frame=none, bgcolor=gray!5]{bash}
tail -f /tmp/arduino_debug.log
\end{minted}

\begin{keybox}
\texttt{LOG\_DEBUG}~--- это \emph{не} \texttt{printf}.
Макрос отключается при сборке с
\texttt{-DNDEBUG}. Так релизная сборка не тратит
время на запись лога.
\end{keybox}

\section{Отладчик \texttt{gdb}}
\label{sec:debug-gdb}

\texttt{gdb}~--- это пошаговый отладчик.
Позволяет остановить программу в любом
месте, посмотреть значения переменных,
пройти по коду строчка за строчкой.

\subsection{Запуск}

Собираем с отладочной информацией:

\begin{minted}[frame=none, bgcolor=gray!5]{bash}
cd arduino-control
mkdir build && cd build
cmake -DCMAKE_BUILD_TYPE=Debug ..
make
\end{minted}

Запускаем под gdb:

\begin{minted}[frame=none, bgcolor=gray!5]{bash}
gdb ./arduino_control
\end{minted}

\subsection{Основные команды}

\begin{table}[htbp]
\centering
\caption{Основные команды \texttt{gdb}}
\label{tab:gdb-commands}
\small
\begin{tabularx}{\textwidth}{@{}l L@{}}
\toprule
\textbf{Команда} & \textbf{Что делает} \\
\midrule
\texttt{run}             & Запустить программу \\
\texttt{break file:line} & Поставить точку останова \\
\texttt{continue}        & Продолжить выполнение \\
\texttt{next}            & Следующая строка (без захода в функции) \\
\texttt{step}            & Следующая строка (с заходом в функции) \\
\texttt{print var}       & Показать значение переменной \\
\texttt{backtrace}       & Показать стек вызовов \\
\texttt{quit}            & Выйти \\
\bottomrule
\end{tabularx}
\end{table}

\subsection{Пример: найти падение}

Проблема: программа падает при открытии
порта. Где именно?

\begin{minted}[frame=none, bgcolor=gray!5]{bash}
$ gdb ./arduino_control
(gdb) run
Program received signal SIGSEGV, Segmentation fault.
0x0000555555558a4c in open_serial_port (app=0x0, port=0x0)
    at src/serial_port.c:245
245     if (!app->is_connected || !app->channel) {
(gdb) backtrace
#0  open_serial_port (app=0x0, port=0x0) at serial_port.c:245
#1  0x0000555555556e10 in on_connect_clicked (button=..., user_data=...)
#2  0x00007ffff7e0a8f5 in g_signal_emit_valist ...
(gdb) print app
$1 = (AppData *) 0x0
\end{minted}

Что видим: \texttt{app == NULL}. Значит,
где-то передали \texttt{NULL} в колбэк.
Ищем источник.

\begin{keybox}
\texttt{backtrace}~--- самая важная команда. Она
показывает \emph{путь} до места падения:
какая функция вызвала какую. Читать
стек~--- базовый навык отладки.
\end{keybox}

\section{Поиск утечек: \texttt{valgrind}}
\label{sec:debug-valgrind}

\texttt{valgrind}~--- это виртуальная машина,
которая перехватывает все вызовы
\texttt{malloc}/\texttt{free} и находит утечки.

\subsection{Запуск}

\begin{minted}[frame=none, bgcolor=gray!5]{bash}
valgrind --leak-check=full \
         --show-leak-kinds=all \
         --track-origins=yes \
         ./arduino_control
\end{minted}

\subsection{Что искать в выводе}

Пример утечки:

\begin{minted}[frame=none, bgcolor=gray!5]{bash}
==12345== HEAP SUMMARY:
==12345==     in use at exit: 1,024 bytes in 8 blocks
==12345==   total heap usage: 12,345 allocs, 12,337 frees,
                567,890 bytes allocated
==12345==
==12345== 128 bytes in 1 blocks are definitely lost in loss record 3 of 8
==12345==    at 0x4C2FB0F: malloc (in /usr/lib/valgrind/...)
==12345==    by 0x108A4C: on_save_log_clicked (callbacks.c:321)
==12345==    by 0x4E3A8F5: g_signal_emit_valist ...
\end{minted}

Что значит:

\begin{itemize}
    \item \textbf{128 bytes}~--- размер утечки.
    \item \textbf{definitely lost}~--- 100\% утечка.
    \item \textbf{on\_save\_log\_clicked (callbacks.c:321)}~---
          место, где выделили и не освободили.
\end{itemize}

\begin{warnbox}
Не все <<утечки>> — настоящие. GLib иногда
аллоцирует память на старте и
не освобождает её при выходе~---
это нормально. Ищите только
\texttt{definitely lost} в своём коде.
\end{warnbox}

\section{Санитайзеры: \texttt{ASan} и \texttt{UBSan}}
\label{sec:debug-sanitizers}

Санитайзеры~--- это встроенные в компилятор
проверки. Они добавляют код, который
ловит ошибки во время выполнения.
Быстрее valgrind~--- в 10-20 раз.

\subsection{Включение}

В нашем \texttt{CMakeLists.txt} они уже включены
для Debug-сборки:

\begin{listing}[htbp]
\begin{minted}[fontsize=\footnotesize, frame=none, bgcolor=gray!5]{cmake}
if(CMAKE_BUILD_TYPE STREQUAL "Debug" OR NOT CMAKE_BUILD_TYPE)
    add_compile_options(-g3 -O0 -fno-omit-frame-pointer)
    add_compile_options(-fsanitize=address,undefined)
    add_link_options(-fsanitize=address,undefined)
endif()
\end{minted}
\caption{Санитайзеры в \texttt{CMakeLists.txt}}
\label{lst:debug-sanitizers}
\end{listing}

\subsection{Что ловит \texttt{ASan}}

\textbf{\texttt{ASan} (AddressSanitizer)} ловит:

\begin{itemize}
    \item \textbf{Переполнение буфера}.
    \item \textbf{Использование после \texttt{free}}.
    \item \textbf{Двойной \texttt{free}}.
    \item \textbf{Утечки (при выходе)}.
\end{itemize}

\textbf{Пример сообщения}:

\begin{minted}[frame=none, bgcolor=gray!5]{bash}
==12345==ERROR: AddressSanitizer: heap-buffer-overflow
READ of size 4 at 0x602000000010 thread T0
    #0 0x4d8b3f in line_buffer_feed line_buffer.c:45
    #1 0x4d7a21 in on_serial_data serial_port.c:112
\end{minted}

Что значит: в \texttt{line\_buffer\_feed} читаем
за пределами буфера. Строка 45~--- место проблемы.

\subsection{Что ловит \texttt{UBSan}}

\textbf{\texttt{UBSan} (UndefinedBehaviorSanitizer)} ловит:

\begin{itemize}
    \item \textbf{Деление на ноль}.
    \item \textbf{Переполнение знаковых чисел}.
    \item \textbf{Сдвиг на отрицательное число}.
    \item \textbf{Выход за границы массива}.
\end{itemize}

\begin{ideabox}
Санитайзеры~--- это <<бесплатные>> проверки. Вы
просто собираете с \texttt{-fsanitize}~---
и получаете сотни проверок во время
работы. Включите их в Debug-сборке~---
и забудьте про большинство ошибок
памяти.
\end{ideabox}

\section{Сообщения \texttt{GLib}}
\label{sec:debug-glib}

\textbf{GLib} пишет предупреждения в
\textbf{stderr. Они выглядят пугающе, но обычно указывают на конкретную проблему.}

\subsection{Типичные сообщения}

\begin{table}[htbp]
\centering
\caption{Типичные предупреждения GLib}
\label{tab:glib-warnings}
\small
\begin{tabularx}{\textwidth}{@{}l L@{}}
\toprule
\textbf{Сообщение} & \textbf{Что значит} \\
\midrule
\texttt{GObject-CRITICAL} & Обращение к уничтоженному объекту \\
\texttt{GLib-GObject-CRITICAL} & Неправильный тип в \texttt{GTK\_*} макросе \\
\texttt{Source ID N was not found} & \texttt{g\_source\_remove} для удалённого таймера \\
\texttt{gtk\_widget\_*: assertion failed} & Виджет не тот, что ожидался \\
\bottomrule
\end{tabularx}
\end{table}

\subsection{Пример: \texttt{GObject-CRITICAL}}

\begin{minted}[frame=none, bgcolor=gray!5]{bash}
(arduino_control:12345): Gtk-CRITICAL :
gtk_widget_set_sensitive: assertion 'GTK_IS_WIDGET (widget)' failed
\end{minted}

Что значит: вызываем \texttt{gtk\_widget\_set\_sensitive}
для объекта, который уже не виджет
(уничтожен). Причина — обращение к UI
после \texttt{on\_window\_destroy}.

Решение — флаг \texttt{ui\_alive}:

\begin{listing}[htbp]
\begin{minted}[frame=none, bgcolor=gray!5]{c}
if (app->ui_alive) {
    gtk_widget_set_sensitive(app->btn_connect, TRUE);
    /* ... */
}
\end{minted}
\caption{Защита через \texttt{ui\_alive}}
\label{lst:debug-ui-alive}
\end{listing}

\begin{keybox}
GLib-CRITICAL~--- это \emph{сигнал}, а не \emph{шум}.
Если видите \texttt{CRITICAL}~--- ищите
причину. Игнорировать их нельзя~---
они предупреждают о реальных багах.
\end{keybox}

\section{Отладка STM32}
\label{sec:debug-stm32}

\textbf{STM32} отлаживается \emph{иначе}: через
\texttt{OpenOCD} и \texttt{gdb}.

\subsection{Подключение}

\textbf{Терминал 1}: \texttt{openocd}.

\begin{minted}[frame=none, bgcolor=gray!5]{bash}
openocd -f interface/stlink.cfg -f target/stm32f4x.cfg
\end{minted}

\textbf{Терминал 2}: \texttt{gdb}.

\begin{minted}[frame=none, bgcolor=gray!5]{bash}
arm-none-eabi-gdb
(gdb) target extended-remote 127.0.0.1:3333
(gdb) file build/firmware.elf
(gdb) load
(gdb) break main
(gdb) continue
\end{minted}

\subsection{HardFault}

\textbf{HardFault}~--- самая частая ошибка на STM32.
\textbf{Что делать}:

\begin{enumerate}
    \item \textbf{Смотрим регистры}:
          \texttt{info registers}.
    \item \textbf{Смотрим PC}~--- где упало.
    \item \textbf{Смотрим CFSR/HFSR}~--- что случилось.
    \item \textbf{Смотрим стек}~--- кто вызвал.
\end{enumerate}

\begin{listing}[htbp]
\begin{minted}[frame=none, bgcolor=gray!5]{bash}
(gdb) info registers
pc             0x08001234
(gdb) x/4z 0x08001234
0x8001234: 0x00000000 0x00000000 ...
(gdb) print/x *(uint32_t*)0xE000ED28
$1 = 0x00008200
\end{minted}
\caption{Диагностика HardFault}
\label{lst:debug-hardfault}
\end{listing}

\begin{keybox}
\textbf{CFSR} (\texttt{0xE000ED28})~--- Configurable
Fault Status Register. \emph{Его} биты \emph{говорят},
\emph{что} \emph{случилось}: \emph{невыровненный}
\emph{доступ}, \emph{деление} \emph{на} \emph{ноль},
\emph{невалидная} \emph{инструкция}.
\end{keybox}

\section{Типичные ошибки}
\label{sec:debug-common}

Собрали список ошибок, которые
встречались в нашем проекте. Каждая~---
с объяснением и решением.

\subsection{Segfault в колбэке}

Симптом: программа падает при клике
на кнопку.

Причина: \texttt{user\_data} равен \texttt{NULL}
или указывает на освобождённую память.

Решение:

\begin{listing}[htbp]
\begin{minted}[frame=none, bgcolor=gray!5]{c}
void on_connect_clicked(GtkButton *button, gpointer user_data)
{
    AppData *app = (AppData *)user_data;
    if (!app) {
        LOG_DEBUG("on_connect_clicked: app == NULL");
        return;
    }
    /* ... */
}
\end{minted}
\caption{Проверка \texttt{user\_data}}
\label{lst:debug-null-check}
\end{listing}

\subsection{Таймер работает после \texttt{app\_data\_free}}

Симптом: \texttt{GObject-CRITICAL} через
несколько секунд после выхода.

Причина: таймер не остановлен до
освобождения \texttt{AppData}.

Решение: \texttt{app\_data\_free} останавливает
все таймеры первым делом.

\subsection{Утечка строк}

Симптом: \texttt{valgrind} показывает
\texttt{definitely lost} в \texttt{gtk\_text\_buffer\_get\_text}.

Причина: \texttt{gtk\_text\_buffer\_get\_text} возвращает
новую строку, которую надо \texttt{g\_free}.

Решение:

\begin{listing}[htbp]
\begin{minted}[frame=none, bgcolor=gray!5]{c}
char *text = gtk_text_buffer_get_text(buffer, &start, &end, FALSE);
/* ... используем text ... */
g_free(text);   /* обязательно */
\end{minted}
\caption{Освобождение строки из \texttt{GtkTextBuffer}}
\label{lst:debug-free-string}
\end{listing}

\subsection{Двойное закрытие \texttt{fd}}

Симптом: \texttt{close(-1)} возвращает \texttt{EBADF}.

Причина: \texttt{close\_serial\_port} вызывается
дважды.

Решение: флаг \texttt{closing\_port} защищает
от рекурсии; проверка \texttt{if (app->fd >= 0)}~---
от повторного закрытия.

\section{Поток отладки}
\label{sec:debug-flow}

Отладка~--- это цикл, а не одноразовое
действие. Схема процесса:

\begin{figure}[htbp]
\centering
\adjustbox{max width=\textwidth}{%
\begin{tikzpicture}[
    x=1cm, y=1cm,
    box/.style={
        draw, rounded corners=2mm, thick,
        minimum width=4.5cm, minimum height=0.9cm,
        align=center, font=\small
    },
    start/.style={box, fill=clrWarn!20},
    action/.style={box, fill=clrLogicData!20},
    check/.style={box, fill=clrIdea!20},
    arr/.style={-{Latex[length=3mm]}, thick}
]

% ─── Координатная сетка (СНАЧАЛА) ───────────────
\draw[step=1cm, gray!20, very thin] (-1, -9) grid (7, 1);

% ─── Блоки ───────────────────────────────────────
\node[start]  (bug)    at (2, 0)    {Проблема};
\node[action] (repro)  at (2, -1.5) {Воспроизвести стабильно};
\node[check]  (hypo)   at (2, -3)   {Гипотеза: <<почему?>>};
\node[action] (check2) at (2, -4.5) {Проверить (gdb, log, valgrind)};
\node[action] (fix)    at (2, -6)   {Исправить};
\node[check]  (verify) at (2, -7.5) {Проверить, что исчезло};

% ─── Стрелки ─────────────────────────────────────
\draw[arr] (bug)    -- (repro);
\draw[arr] (repro)  -- (hypo);
\draw[arr] (hypo)   -- (check2);
\draw[arr] (check2) -- (fix);
\draw[arr] (fix)    -- (verify);
\draw[arr] (verify.east) -- ++(2, 0) |- (hypo.east);
\node[draw=none, font=\footnotesize, right] at (3.6, -5.3)
    {если не помогло};

% ─── Метки осей (БЕЗ рамок) ─────────────────────
\foreach \x in {-1,...,7} {
    \node[draw=none, font=\tiny, gray] at (\x, -9.3) {\x};
}
\foreach \y in {-9,...,1} {
    \node[draw=none, font=\tiny, gray] at (-1.3, \y) {\y};
}

\end{tikzpicture}
}% end adjustbox
\caption{Цикл отладки.}
\label{fig:debug-flow}
\end{figure}


\begin{ideabox}
<<Воспроизвести стабильно>>~--- \emph{первый} шаг.
Если баг не воспроизводится~--- нельзя
проверить, что он исправлен. Плавающие
баги лечатся не отладчиком, а
логированием~--- ищите закономерность.
\end{ideabox}

\section{Уроки главы}
\label{sec:debug-lessons}

\begin{enumerate}
    \item Отладка~--- это \emph{навык}, а не магия.
          Тренируется, как и любой другой.
    \item \texttt{debug\_log}~--- первый инструмент.
          Работает всегда, не требует внешних программ.
    \item \texttt{backtrace}~--- самая важная команда
          \texttt{gdb}. Указывает путь до падения.
    \item \texttt{valgrind}~--- для утечек. Ищите
          только \texttt{definitely lost} в своём коде.
    \item Санитайзеры~--- бесплатные проверки.
          Включайте в Debug-сборке.
    \item \texttt{GLib-CRITICAL}~--- сигнал, а не шум.
          Ищите причину.
    \item <<Воспроизвести стабильно>>~--- первый шаг.
          Плавающие баги лечатся логированием.
    \item \textbf{STM32} отлаживается через
          \texttt{OpenOCD} + \texttt{gdb}.
\end{enumerate}

\section{Чек-лист грабель}
\label{sec:debug-pitfalls}

\begin{table}[htbp]
\centering
\caption{Грабли отладки}
\label{tab:debug-pitfalls}
\small
\begin{tabularx}{\textwidth}{@{}c L L@{}}
\toprule
\textbf{№} & \textbf{Грабля} & \textbf{Решение} \\
\midrule
1 & \texttt{segfault} без сообщения & \texttt{gdb} + \texttt{backtrace} \\
2 & Утечка памяти & \texttt{valgrind --leak-check=full} \\
3 & Плавающий баг & \texttt{LOG\_DEBUG} + \texttt{tail -f} \\
4 & \texttt{GObject-CRITICAL} в логе & Проверить \texttt{ui\_alive} \\
5 & \texttt{Source ID not found} & Не удалять дважды \texttt{g\_source\_remove} \\
6 & \texttt{fd = -1} vs \texttt{fd = 0} & Явно инициализировать в \texttt{app\_data\_init} \\
7 & \texttt{close(-1)} & Проверять \texttt{fd >= 0} \\
8 & Утечка из \texttt{gtk\_text\_buffer\_get\_text} & \texttt{g\_free(text)} \\
9 & \texttt{HardFault} на STM32 & \texttt{CFSR}, \texttt{HFSR}, \texttt{backtrace} \\
10 & Не видно \texttt{printf} из прошивки & \texttt{SWO} или \texttt{UART} \\
\bottomrule
\end{tabularx}
\end{table}

\section{Упражнения}
\label{sec:debug-exercises}

\begin{exercisebox}
\textbf{Упражнение 1.} Соберите проект с \texttt{-fsanitize=address}
и запустите. Найдите все утечки. Сколько их в \emph{вашем}
коде (без учёта GLib и GTK)?
\end{exercisebox}

\begin{exercisebox}
\textbf{Упражнение 2.} Поставьте точку останова в
\texttt{on\_serial\_line\_received} и запустите программу
под \texttt{gdb}. Понаблюдайте, что происходит при приёме
строки. Сколько \emph{уровней} \texttt{backtrace} до
\texttt{on\_serial\_data}?
\end{exercisebox}

\needspace{5\baselineskip}
\begin{exercisebox}
\textbf{Упражнение 3.} Добавьте в \texttt{debug\_log} поддержку
\emph{уровней}: \texttt{DEBUG}, \texttt{INFO}, \texttt{ERROR}.
Как фильтровать сообщения по уровню через \emph{переменную
окружения}?
\end{exercisebox}

\begin{exercisebox}
\textbf{Упражнение 4.} Что произойдёт, если вызвать
\texttt{close\_serial\_port} из \emph{двух} потоков
одновременно? Как защитить функцию через \texttt{GMutex}?
\end{exercisebox}

\begin{exercisebox}
\textbf{Упражнение 5.} Напишите bash-скрипт, который
запускает программу под \texttt{valgrind} и \emph{парсит}
вывод~--- выводит только \texttt{definitely lost} с указанием
файла и строки. Как автоматизировать поиск утечек?
\end{exercisebox}


